% from ega4-13-fr.tex
\oldpage[IV-3]{195}
  \item Il existe un entier $e$ et un voisinage ouvert $U$ de $x$ tels que
  l'image par $f$ de toute composante irréductible de $U$ soit dense dans un
  $Y_\alpha$, et que, pour tout $x'\in U$, l'espace
  $U\cap f^{-1}(f(x'))$ soit équidimensionnel et de dimension $e$.
  \item[\emph{a$'$)}] Il existe un entier $e$ et un voisinage ouvert $U$ de
  $x$ tels que l'image par $f$ de toute composante irréductible de $U$ soit
  dense dans un $Y_\alpha$ et que, si on désigne par $y_\alpha$ le point
  générique de $Y_\alpha$, toute composante irréductible des espaces
  $U\cap f^{-1}(y)$, $U\cap f^{-1}(y_\alpha)$ soit de dimension $e$.
  \item[\emph{a$''$)}] Il existe un entier $e$ et un voisinage ouvert $U$ de
  $x$ tels que, pour chacune des composantes irréductibles $U_\lambda$ de
  $U$, $f(U_\lambda)$ soit dense dans un $Y_\alpha$ et que, pour tout
  $x'\in U_\lambda$, les composantes irréductibles de
  $U_\lambda\cap f^{-1}(f(x'))$ soient toutes de dimension $e$.
  \item Il existe un entier $e$, un voisinage ouvert $U$ de $x$ et un
  $Y$-morphisme quasi-fini
  \[
    g:U\longrightarrow Y\otimes_{\mathbf Z}\mathbf Z[T_1,\ldots,T_e]
  \]
  (préschéma que nous noterons aussi $Y[T_1,\ldots,T_e]$ pour abréger) tels
  que l'image par $g$ de toute composante irréductible de $U$ soit dense dans
  une composante irréductible de $Y[T_1,\ldots,T_e]$.
\end{enumerate}
\end{proposition}

\begin{proof}
Il est immédiat que $a''$) entraîne $a$), car pour tout $x\in U$, les
composantes irréductibles de $U\cap f^{-1}(f(x))$ sont chacune composante
irréductible d'un des $U_\lambda\cap f^{-1}(f(x))$, d'où la conclusion par
($\mathbf 0$, 14.1.4). La condition $a$) entraîne trivialement $a'$).
Montrons ensuite que $a'$) entraîne $a''$)~; on peut se borner au cas où $f$
est de type fini et $X$ et $Y$ réduits (1.5.4). Soit $U_\lambda$ une
composante irréductible de $U$, et supposons que $f(U_\lambda)$ soit dense
dans $Y_\alpha$~; alors la restriction de $f$ à $U_\lambda$ se factorise en
$U_\lambda\xrightarrow{f_\alpha}Y_\alpha\to Y$, où $f_\alpha$ est de type
fini et dominant (\textbf{I}, 5.2.2). Soit $x_\lambda$ le point générique de
$U_\lambda$~; en vertu de ($\mathbf 0_{\mathrm I}$, 2.1.8),
$U_\lambda\cap f_\alpha^{-1}(y_\alpha)$ est l'unique composante irréductible
de $U\cap f^{-1}(y_\alpha)$ contenant $x_\lambda$ et est par hypothèse de
dimension $e$, égale aux dimensions de toutes les composantes irréductibles
de $U_\lambda\cap f^{-1}(y)$, en vertu de l'hypothèse et de (13.1.1). Mais en
vertu du théorème de Chevalley (13.1.3) et de (13.1.1), l'ensemble $V_\lambda$
des $x'\in U_\lambda$ tels que
$\dim_{x'}(f_\lambda^{-1}(f_\lambda(x')))=e$ est ouvert et contient $x$, et il
suffit de prendre la réunion des $V_\lambda$ pour obtenir un ensemble ouvert
satisfaisant aux conditions de $a''$).

Prouvons maintenant que $a$) entraîne $b$)~; on peut se limiter au cas où
$Y=\operatorname{Spec}(A)$ et $X=\operatorname{Spec}(B)$ sont affines et où
$U=X$. Prouvons d'abord le lemme suivant~:

\begin{lemma}[13.3.1.1]
\label{IV.13.3.1.1-fr}
Soient $Y=\operatorname{Spec}(A)$, $X=\operatorname{Spec}(B)$ deux schémas
affines, $f:X\to Y$ un morphisme de type fini, $y$ un point de $f(X)$, et
posons $e=\dim(f^{-1}(y))$. Alors il existe un $Y$-morphisme
$g:X\to Y[T_1,\ldots,T_e]$ tel que (si l'on pose
$X_y=X\times_Y\operatorname{Spec}(k(y))$), le morphisme
$g_y:X_y\to\operatorname{Spec}(k(y)[T_1,\ldots,T_e])$ soit fini. En outre,
pour un tel morphisme $g$, $g_y$ est nécessairement surjectif et il existe un
voisinage ouvert $U$ de $f^{-1}(y)$ dans $X$ tel que $g|U$ soit quasi-fini.
\end{lemma}

Soit $\mathfrak p=\mathfrak j_y$~; l'anneau $B\otimes_Ak(\mathfrak p)$ est une
$k(\mathfrak p)$-algèbre de type fini, donc le lemme de normalisation
(Bourbaki, \emph{Alg. comm.}, chap.~V, \S~3, n\textsuperscript{o}~1, th.~1)
prouve qu'il y a dans $B\otimes_Ak(\mathfrak p)$ une suite finie
$(t_i)_{1\leq i\leq r}$ d'éléments algébriquement indépendants sur
$k(\mathfrak p)$ et tels que si l'on pose
$C'=k(\mathfrak p)[t_1,\ldots,t_r]$, $B\otimes_Ak(\mathfrak p)$ soit une
$C'$-algèbre \emph{finie}~; on a donc
$\dim(B\otimes_Ak(\mathfrak p))=\dim(C')$ ($\mathbf 0$, 16.1.5) et comme
$\dim(C')=r$ (5.2.1), on a $r=e$. Comme
$B\otimes_Ak(\mathfrak p)=(B/\mathfrak pB)\otimes_{A/\mathfrak p}k(\mathfrak p)$,
on peut, en multipliant les $t_i$ par un élément $\neq0$

